\chapter*{ЗАКЛЮЧЕНИЕ}

Исследование показало, что гибридные силовые установки являются эффективным инструментом повышения топливной экономичности автомобилей, позволяя совмещать преимущества ДВС и электродвигателей. Было установлено, что внедрение электропривода приводит к значительному снижению расхода топлива и выбросов в атмосферу, особенно в городских условиях, где рекуперация энергии при торможении и возможность движения на электротяге оказываются наиболее востребованными.

По итогам классификации гибридных силовых установок были выявлены различия между их типами и определены степени гибридизации. Был проведён анализ применения и перспектив развития различных классов гибридов, который показал тенденцию рынка в сторону подключаемых и мягких гибридных систем. Производители постепенно сокращают долю традиционных автомобилей с ДВС, при этом роль полных гибридов в модельных линейках также снижается.

В ходе работы была изучена современная терминология в области гибридных силовых установок, раскрыты понятия последовательной, параллельной и комбинированной схем, а также рассмотрены области их применения. Были проанализированы основные архитектуры гибридных трансмиссий, выявлены их конструктивные особенности и принципы работы.

Далее был проведён анализ эффективности различных архитектур гибридных силовых установок на основе собранных данных по расходу топлива. Было установлено, что преимущества гибридных систем особенно заметны в городском цикле движения. Наибольшую экономию топлива демонстрируют архитектуры HSD и Multi-mode Series, тогда как Parallel Inline обеспечивает стабильную эффективность на шоссе. MHEV показывают умеренное улучшение экономичности, но при этом значительно повышают динамические характеристики автомобиля.

Таким образом, было показано, что гибридные системы предоставляют широкий спектр решений для повышения экономичности и экологичности автомобилей, позволяя адаптировать конструкции под различные условия эксплуатации. Результаты исследования подтвердили, что архитектура трансмиссии и степень интеграции электродвигателей являются ключевыми факторами, определяющими эффективность гибридного автомобиля.

\addcontentsline{toc}{chapter}{ЗАКЛЮЧЕНИЕ}